% Part: lambda-calculus
% Chapter: introduction
% Section: overview

\documentclass[../../../include/open-logic-section]{subfiles}

\begin{document}

\olfileid{lam}{int}{ovr}
\olsection{لنډه کتنه}

الونزو چرچ په لومړيو ۱۹۳۰مو کلونو کښې لامبډا حساب په اصل کښې د
تعميري منطق د بنسټ په توګه جوړ کړے و، \emph{نۀ} د محاسبه کېدونکو
تابعو د نمونې په توګه۔ خو ډېر ژر وښودل شوه چې دا د محاسبه کېدنې له
نورو تعريفونو، لکه د ټيورينګ له مخې محاسبه کېدونکو تابعو او جزوي
بازګشتي تابعو، سره معادل دے۔ دا چې دغه معادلتوب په لومړي سر کښې يو
څه حيرانوونکے و، د دې ځانګړنې ارزښت لا زياتوي۔

د لامبډا نښه د تابعې د نوم له جلا تعريفولو پرته، د هغې په تعريفوونکي
نښيز عبارت سره نېغې تابعې ته د اشارې يوه آسانه لاره ده۔ د دې پر ځاے
چې ووايو ``$f$ هغه تابع وګڼئ چې $f(x) = x + 3$ ئې تعريفوي''، ويلے
شو ``$f$ هغه تابع وګڼئ چې $\lambd[x][(x + 3)]$ ئې تعريفوي''۔ په بله
وينا، $\lambd[x][(x+3)]$ يوازې د هغې تابعې يو \emph{نوم} دے چې په
خپل آرګومېنټ درې ورزياتوي۔ په دې عبارت کښې $x$ يو تش‌ځايے متغير دے:
هماغه تابع په $\lambd[y][(y + 3)]$ هم ښودل کېدے شي۔ دا نښه د نورو
پاراميټرونو په شتون کښې هم کار کوي۔ مثلاً، که $g(x, y)$ يوه دوه‌ځاييزه
تابع او $k$ يو طبيعي عدد وي، نو $\lambd[x][g(x,k)]$ هغه تابع ده چې
هر~$x$،~$g(x, k)$ ته وړي۔

له نښيز عبارت څخه د تابعې د تعريف دغه طريقه \emph{لامبډا تجريد}
بلل کېږي۔ د لامبډا تجريد بل اړخ \emph{تطبيق} دے: که يوه تابع~$f$
ولرو (مثلاً پر طبيعي عددونو تعريف شوې)، پر هر قيمت، لکه ۲، ئې تطبيقولے
شو۔ په دوديزه نښه کښې د پايلې دپاره البته~$f(2)$ ليکو۔

که لامبډا تجريد او تطبيق سره يوځاے کړو، څۀ پېښېږي؟ بيا حاصل عبارت
په دې ساده کېدے شي چې د تجريد شوي متغير په ځاے د تابعې ورکړل شوے
آرګومېنټ کېږدو۔ د بېلګې په توګه،
\[
(\lambd[x][(x + 3)])(2)
\]
تر~$2 + 3$ پورې ساده کېدے شي۔

تر دې ځايه مو دوديزو مفاهيمو ته يوازې نوې نښې ورکړې دي۔ خو لامبډا
حساب د سټونو د تيورۍ له ليدلوري څخه ډېر بنسټيز بېلتون لري۔ په دې
چوکاټ کښې:
\begin{enumerate}
\item هر څۀ يوه تابع ښيي۔
\item تابعې د لامبډا تجريد په وسيله تعريفېدے شي۔
\item هر څۀ پر هر بل څيز تطبيقېدے شي۔
\end{enumerate}
مثلاً، که $F$ د لامبډا حساب يو ترم وي، $F(F)$ تل د معنٰی لرونکے
ګڼل کېږي۔ دې پراخ چوکاټ ته \emph{بې‌ټايپه} لامبډا حساب وايي؛ دلته
``بې‌ټايپه'' معنا دا ده چې پر څۀ د څۀ شي د تطبيقولو محدوديت نشته۔

\begin{digress}
يو \emph{ټايپ‌لرونکے} لامبډا حساب هم شته چې د بې‌ټايپې بڼې يو مهم
بدلون دے۔ که څۀ هم ټايپ‌لرونکے لامبډا حساب په ډېرو اړخونو کښې
بې‌ټايپې بڼې ته ورته دے، له کلاسیک سټي چوکاټ سره ئې سمون ډېر آسان
دے او ځينې خاصيتونه ئې ډېر توپير لري۔

د لامبډا حساب څېړنه د نظري کمپيوټر ساينس او د پروګرامي ژبو د طرحې
په منځ کښې مهمه ثابته شوې ده۔ LISP، چې جان مک‌کارتي په ۱۹۵۰مو
کلونو کښې جوړه کړې وه، د داسې مفکورو تر اغېز لاندې د يوې لومړنۍ
پروګرامي ژبې بېلګه ده۔
\end{digress}

\end{document}
